\documentclass{article}
\usepackage[a4paper, total={6.5in, 10in}]{geometry}
\setlength{\parindent}{0pt}
\usepackage{tcolorbox}
\tcbset{colback=white!94!black, colframe=black!60!white, coltitle = white, boxrule=0.8pt, arc=2mm}
\newtcolorbox{boxs}[2][]{title= #2,#1}
\usepackage{datetime2}
\usepackage[british]{babel}
\usepackage{amsfonts}
\usepackage{amsmath}

\title{\textbf{Linear Algebra Homework }}
\author{Robert O'Connell}
\date{\today}
\begin{document}
\maketitle
\vspace{5mm}
\section*{Question 1}
\[
\vec{u} = \langle 1,4,-7 \rangle, \quad \vec{v} = \langle 2,-1,4 \rangle, \quad \vec{w} = \langle 1,5,-1 \rangle, \quad \vec{b} = \langle 0,-9,18 \rangle
\]
\subsubsection*{(a) Computing the area of the parallelogram determined by the vectors \(\vec{u}\) and \(\vec{v}\)}
The area determined by two vectors is given by the magnitude of their cross product
\begin{align*}
    ||\vec{u} \times \vec{v}|| &= ||\langle u_2v_3 - u_3v_2, u_3v_1 - u_1v_3, u_1v_2 - u_2v_1 \rangle|| \\
    ||\vec{u} \times \vec{v}|| &= ||\langle 16 - 7, -14 - 4, -1 -8 \rangle||\\
    ||\vec{u} \times \vec{v}|| &= ||\langle 9, -18, -9 \rangle||\\
    ||\vec{u} \times \vec{v}|| &= \sqrt{(9)^2 + (-18)^2 + (-9)^2}\\
    ||\vec{u} \times \vec{v}|| &= \sqrt{486}\\
\end{align*}

\subsubsection*{(b) Computing the volume of the parallelepiped determined by the vectors \(\vec{u}\), \(\vec{v}\) and \(\vec{w}\)}
The volume of a parallelepiped determined by three vectors is the absolute value of the cross product of two vectors dotted with the third vector
\begin{align*}
    |(\vec{u} \times \vec{v}) \cdot \vec{w}| &= |\langle 9, -18, -9 \rangle \cdot \langle 1,5,-1 \rangle | \\
    |(\vec{u} \times \vec{v}) \cdot \vec{w}| &= | 9 - 90 + 9| \\
    |(\vec{u} \times \vec{v}) \cdot \vec{w}| &= 72 
\end{align*}

\subsubsection*{(c) Finding     \(|\vec{b} \cdot (\vec{u} \times \vec{v})|\)}
\begin{align*}
    |\vec{b} \cdot (\vec{u} \times \vec{v})| &= |\langle 0,-9,18 \rangle \cdot \langle 9, -18, -9 \rangle| \\
    &= 0 + 162 - 162 \\
    &= 0
\end{align*}
Thus, we can conclude that the vectors \(\vec{u}, \vec{v}\) and \(\vec{b}\) are coplanar

\section*{Question 2}
\subsubsection*{Finding the vector equation of a line given a point and a parallel line}
We are given a parallel line so finding a vector between any two points on the line is the direction
vector for the other line 

Getting the vector between \(t = 0\) and \(t = 1\) for the parallel line
\begin{align*}
    \vec{v} &= \langle -1 - 2, 3 + 1, 7 - 8 \rangle \\
    &= \langle -3,4,-1 \rangle
\end{align*}
We now have both a point on the line and a direction vector and thus are able to find the equation of the line
\begin{align*}
    \vec{r} &= \langle -1,2,4 \rangle + t\langle -3,4,-1 \rangle \\
    &= \langle -3t-1, 4t+2, 4-t \rangle
\end{align*}

\section*{Question 3}
\subsubsection*{Finding the equation of a plane given a point and a perpendicular line}
A plane is described by a point in the plane and a normal vector

We are given a line perpendicular to the plane, and so a direction vector of the line is the normal 
vector of the plane

Take the direction vector between \(t = 0\) and \(t = 1\) as our normal vector
\begin{align*}
    \vec{n} &= \langle -3, 2 -1, 3-7 \rangle \\
    &= \langle -3,1,-4 \rangle
\end{align*}
Let \(P(x,y,z)\) be any other point on the plane and let \(P_0(2,0,1)\) be the known point. Then,
\[
\overrightarrow{P_0P} \cdot \vec{n} = 0
\]

\[
\langle x -2,y,z-1 \rangle \cdot \langle -3,1,-4 \rangle = 0
\]
\[
6-3x + y + 4 -4z = 0
\]
\[
3x -y +4z = 10
\]

Thus, \(3x -y +4z = 10\) is the equation of the plane

\section*{Question 4}
\subsubsection*{Finding the equation of a plane which passes through three given points}
\[
P_1 =(1,1,0), \quad P_2 = (2,0,1),\quad P_3 = (0,-1,3)
\]
To describe the plane, we need a normal vector, which we can obtain by getting the cross product of two vectors
in the plane

\[
\overrightarrow{P_1P_2} = \langle 1,-1,1 \rangle, \quad \overrightarrow{P_2P_3} = \langle -2,-1,2 \rangle
\]

\begin{align*}
    \vec{n} &= \overrightarrow{P_1P_2} \times \overrightarrow{P_2P_3} \\
    &= \langle -2 +1,-2 -2, -1-2 \rangle \\
    &= \langle -1,-4,-3 \rangle
\end{align*}

Let \(P(x,y,z)\) be any other point in the plane
\[
\overrightarrow{PP_1} \cdot \vec{n} = 0
\]
\[
\langle x-1,y-1,z \rangle \cdot \langle -1,-4,-3 \rangle = 0
\]
\[
1-x,4-4y, -3z = 0
\]
\[
x+4y+3z = 5
\]
Thus, \(x+4y+3z = 5\) is the equation of the plane

\section*{Question 5}
\subsubsection*{Finding the intersection of a line and a plane}
Let the vector between the two points \((\vec{v})\) be the direction vector
\[
\vec{v} = \langle 3,-2,1 \rangle
\]
Then the equation of the line is
\begin{align*}
    \vec{r} &= \langle 1,0,1 \rangle + t\langle 3,-2,1 \rangle \\
    &= \langle 1+3t, -2t, 1+t \rangle
\end{align*}
The parametric equations are
\[
x = 1+3t 
\]
\[
y= -2t
\]
\[
z = 1+t
\]

Substituting these into the equation of the plane, to find the point of intersection
\[
x +y+z = 6
\]
\[
1+3t-2t+1+t = 6
\]
\[
t = 2
\]
The point at \(t = 2\) is \((7, -4, 3)\), which is the point of intersection

\section*{Question 6}
\subsubsection*{Do three given points form a right triangle?}
\[
P_1 = (-1,7,9), \quad P_2 = (1,5,14),  \quad P_3 = (-9,11,13)
\]

Any three points in \(\mathbb{R}^3\) are coplanar, and such a triangle is formed between them

So the only thing to check is whether or not there is a right angle

Check this by checking if any two of the vectors between the three points are perpendicular to each other
\[
\overrightarrow{P_1P_2} = \langle 2,-2,5 \rangle, \quad \overrightarrow{P_2P_3} = \langle -10,6,-1 \rangle, 
\quad \overrightarrow{P_1P_3} = \langle -8,4,4 \rangle
\]

\begin{align*}
    \overrightarrow{P_1P_2} \cdot \overrightarrow{P_2P_3} &= -20-12-5 \\
    &= -37 \\
    &\neq 0
\end{align*}
\begin{align*}
    \overrightarrow{P_1P_2} \cdot \overrightarrow{P_1P_3} &= -16-8+20 \\
    &= -4 \\
    &\neq 0
\end{align*}
\begin{align*}
    \overrightarrow{P_1P_3} \cdot \overrightarrow{P_2P_3} &= -80 +24-4 \\
    &= -60 \\
    &\neq 0
\end{align*}

Since none of the three vectors are perpendicular to one another, we conclude that a right triangle is not formed

\section*{Question 7}
\subsubsection*{Find three points that do form a right triangle in \(\mathbb{R}^3\)}

Let,
\[
P_1 = (-1,7,9), \quad P_2 = (1,5,14),  \quad P_3 = (x,y,z)
\]

Then,
\[
\overrightarrow{P_1P_2} = \langle 2,-2,5 \rangle, \quad \overrightarrow{P_2P_3} = \langle x-1,y-5,z-14 \rangle, 
\quad \overrightarrow{P_1P_3} = \langle x+1,y-7,z-9 \rangle
\]
For there to be a right angle any two vectors should be perpendicular

\[
\overrightarrow{P_1P_2} \cdot \overrightarrow{P_2P_3} = 0
\]

\[
2x-2-2y+10+5z-70 = 0 
\]
\[
2x-2y+5z = 62
\]

For the three points to form a right triangle, \(P_3\) needs to be in the plane \(2x-2y+5z = 62\)

One such point is 
\[P_3 =(31,0,0)\]

Thus the points \((-1,7,9), (1,5,14)\) and \((31,0,0)\) form a right triangle 
\end{document}